\subsection{An elementary ideal-lattice decay lemma}
\label{app:ag-lattice}

This argument is a derived alternative to the archive's
lattice-sum step; it uses only the algebraic norm and elementary
packing.  It does not prove the representation-theoretic Fourier
estimates to which the lemma is applied.

\paragraph{Statement.}
Let \(E\) be a totally complex number field of degree \(2f\),
\(f\ge1\), and choose exactly one embedding
\(\sigma_\nu:E\hookrightarrow\mathbb C\) from each complex-conjugate
pair.  Let \(\mathfrak a\ne0\) be a fractional ideal.  Its absolute
norm \(N\mathfrak a\) is the positive rational norm, extended
multiplicatively from integral ideals.  For \(x_\nu>0\), \(K\ge3\)
real, and \(\Theta\ge1\), assume
\[
 \prod_{\nu=1}^{f}x_\nu\le
        \frac{\sqrt{N\mathfrak a}}{\Theta^f}.
\]
Then
\[
 \sum_{\substack{z\in\mathfrak a\\z\ne0}}
 \prod_{\nu=1}^{f}
       \min\{1,(x_\nu/|\sigma_\nu z|)^K\}
 \le \bigl(384\,\Theta^{-(K-1/2)}\bigr)^f.                 \tag{AG.A1}
\]
The bound is formally valid at \(\Theta=1\); its smallness is
useful only when \(\Theta\) is sufficiently large compared with
the fixed constant.

\paragraph{Proof.}
For \(0\ne z\in\mathfrak a\), the ideal
\((z)\mathfrak a^{-1}\) is integral.  Its norm is a positive
integer, so
\[
 |N_{E/\mathbb Q}z|\ge N\mathfrak a,\qquad
 \prod_{\nu=1}^{f}|\sigma_\nu z|\ge\sqrt{N\mathfrak a}.
\]
The square root is essential: each chosen complex embedding
represents two embeddings in the field norm.
Write \(Z_\nu(z)=\sigma_\nu(z)/x_\nu\).
For distinct \(z,z'\in\mathfrak a\),
\[
 \prod_\nu|Z_\nu(z)-Z_\nu(z')|\ge\Theta^f.                 \tag{AG.A2}
\]

For radii \(R_\nu\ge1\), put \(P=\prod_\nu R_\nu\) and let
\(M(R)\) count the nonzero \(z\in\mathfrak a\) with
\(|Z_\nu(z)|\le R_\nu\) for every \(\nu\).
If \(P<\Theta^f\), then \(M(R)=0\).
Otherwise put \(q=P^{1/f}/\Theta\ge1\) and
\(h=\lceil2q\rceil\).
At each complex place partition the square
\([-R_\nu,R_\nu]^2\) into half-open squares of side \(R_\nu/h\),
assigning the outer boundaries consistently.
At most \((2h+2)^2\le(8q)^2\) squares are needed.
Two points lying in the same product of these squares would
satisfy
\[
 \prod_\nu|Z_\nu(z)-Z_\nu(z')|
 \le(\sqrt2/h)^fP
 \le 2^{-f/2}\Theta^f<\Theta^f,
\]
contradicting (AG.A2).  Thus each product box contains at most
one point, and the packing is in \(2f\) real dimensions:
\[
 M(R)\le64^f\,P^2/\Theta^{2f}
       \quad\hbox{if }P\ge\Theta^f.                        \tag{AG.A3}
\]

First take \(K=3\).  The nonnegative layer-cake identity
\[
 \min\{1,t^{-3}\}=\int_1^\infty 3R^{-4}\,
                                      \mathbf1_{t\le R}\,dR
 \qquad(t>0)
\]
and Tonelli's theorem give, using (AG.A3),
\begin{align*}
 \sum_{z\ne0}\prod_\nu\min\{1,|Z_\nu(z)|^{-3}\}
 &\le 192^f\Theta^{-2f}
 \int_{\substack{R_\nu\ge1\\\prod R_\nu\ge\Theta^f}}
                         \prod_\nu R_\nu^{-2}\,dR_\nu\\
 &=192^f\Theta^{-2f}
 \int_{\substack{t_\nu\ge0\\\sum t_\nu\ge f\log\Theta}}
                            e^{-\sum t_\nu}\,dt\\
 &\le192^f\Theta^{-5f/2}
       \int_{t_\nu\ge0}e^{-\frac12\sum t_\nu}\,dt\\
 &=384^f\Theta^{-5f/2}.
\end{align*}
The substitution in the second line is \(R_\nu=e^{t_\nu}\).
For general \(K\ge3\), every nonzero lattice point has
\(\prod_\nu\max\{1,|Z_\nu(z)|\}\ge\Theta^f\).  Consequently
\[
 \prod_\nu\min\{1,|Z_\nu(z)|^{-K}\}
 \le \Theta^{-(K-3)f}
                   \prod_\nu\min\{1,|Z_\nu(z)|^{-3}\}.
\]
Multiplying the preceding bound proves (AG.A1).

\paragraph{Relation to the archived estimate.}
The source's \texttt{latSum\_of\_core} proves the sufficient
bound \((C\Theta^{-(K-5/2)})^f\), for integer \(K\ge3\)
and \(\Theta>2\).  Formula (AG.A1) implies this form directly
with \(C=384\).  Thus it can replace that single lattice-sum
step in the human argument while all downstream estimates are
left at their archived strength.  Its validity does not depend
on the root discriminant, on a chosen integral basis, or on the
shape of the embedded fractional ideal.

\subsection{The precise remaining local spectral input}
\label{app:ag-local}

To avoid concealing the difficult part inside the phrase
\emph{Fourier decay}, we record the scales used by the source.
For each place write \(I_\nu=[a_\nu,b_\nu]\) and
\(I_{\nu,1}=[a_{\nu,1},b_{\nu,1}]\subset I_\nu\), put
\(N=20(r+1)\), and let
\[
 B_\nu(Y)=\kappa_\nu^{-(2r+1)}
 \sqrt{\left|
 \frac{P_{\lambda_\nu}(Y)P_{\lambda_\nu}(Y+1)}
 {-T_uT_v\prod_{j=1}^k
       (Y+\tfrac12+u_{\nu j})(Y+\tfrac12+v_{\nu j})}
                    \right|}.
\]
Here \(\kappa_\nu=2\pi/|\sigma_\nu(g_0g_1)|\);
it is not the global lattice covolume \(\kappa\) in (AG.4).
The three types of local sizes are
\begin{align*}
 X_{i\nu}
 &=C\sup_{Y\in I_{\nu,i}}B_\nu(Y)
       \left(\frac{C\kappa_\nu N}{b_\nu-a_\nu}\right)^{2i-2},
 & I_{\nu,i}&=\begin{cases}I_{\nu,1}&i=1,\\I_\nu&i>1,\end{cases}\\
 \Lambda_{i,s,\nu}
 &=C\sup_{Y\in I_\nu}B_\nu(Y)
       \left(\frac{C\kappa_\nu N}{b_\nu-a_\nu}\right)^{2(i-1)+s-1},\\
 \Lambda'_{i\nu}
 &=\left[
 C\sup_{Y\in I_\nu}B_\nu(Y)\,
 \kappa_\nu^{2i-1}(T_uT_v)^{-1/2}
       \left(\frac{CN}{b_\nu-a_\nu}\right)^{2(i-1)}
        \right]^{1/2}.
\end{align*}
The source defines \(\Theta^{-f}\) as the maximum of
\[
 \frac{\prod_\nu X_{i\nu}}{\sqrt{N\mathfrak a^V_{i,1}}},
 \qquad
 \frac{\prod_\nu\Lambda_{i,s,\nu}}
 {\min\{\sqrt{N\mathfrak a^V_{i,s}},
                   \sqrt{N\mathfrak a^P_{i,s-1}}\}},
 \qquad
 \frac{\prod_\nu\Lambda'_{i\nu}}
                   {\sqrt{N\mathfrak a^P_{i,r+2-i}}},
                                                               \tag{AG.A4}
\]
over respectively \(1\le i\le r+1\);
\(i\ge1,s\ge2,i-1+s\le r+1\);
and \(1\le i\le r\).
The ideals contain the appropriate Fourier-dual or
inverse-difference lattices.

The local theorem assumes interlacing of both compact
parameters with \(\lambda_\nu\), a cut as in (AG.5), a common
interlacing parameter one rank lower, and
\[
 \begin{split}
 b_\nu-a_\nu&\ge
 C[kN(r+1)^2+N^2(r+1)^2+N^3],\\
 b_{\nu,1}-a_{\nu,1}&\ge N^3/2,\\
 \operatorname{gap}_{\rm nc}(\lambda_\nu)&\ge
 C\,m^{9/2}\bigl(NR^2+R^3\bigr),
 \end{split}                                                   \tag{AG.A5}
\]
where \(m=2r+2\), \(R\) denotes the source's total rank, and
the noncompact gap is the minimum of
\(|\lambda_{\nu i}-\lambda_{\nu j}|\) over pairs whose signature
signs differ.
It also uses the compatible discrete-series and compact-type
representations, smooth finite-type coefficients, positive
intermediate forms, and a compatible smooth big-cell section.
With at most \(N=4K(r+1)\) letters available, it proves:

\begin{enumerate}
 \item the zero Fourier mode and the decay estimate (AG.6);
 \item an off-diagonal bound with profile
 \(\prod_\nu\min\{1,
   (\Lambda_{i-s+1,s+1,\nu}/|t_\nu^{-1}-(t'_\nu)^{-1}|)^K\}\)
 below the top level, and \(2\Lambda'_{i-s+1,\nu}\) in place
 of \(\Lambda\) at the top;
 \item a block-of-size-\(s\) lattice-sum bound with factor
 \(C^{fs}\) times the product of the preceding spectral
 volumes, and a probability average of extensions by at
 most \((s+1)K\) letters.
\end{enumerate}

The off-diagonal estimate has at most \(3^f\) possible place
patterns, absorbed into the fixed constant raised to \(f\).
The probability distributions are independent of the finite
point set and of its positive coefficient kernel.  That
uniformity is needed for the final row estimate.
The archive derives these assertions through Capelli and
raising/lowering identities, positivity, interpolation and
integration by parts.  The present reconstruction has read their
precise statements and the subsequent Poisson/recurrence proof;
it does not yet supply a complete human proof of these local
representation-theoretic assertions.
